\section{Onde aprofundar}

Os títulos marcados com \pt têm edição em português.

\subsection*{Livros}
\small
\begin{longtable}{@{}p{1.8cm} p{4.4cm} p{2.6cm} p{4cm}@{}}
\toprule
Tema & Livro & Autor(es), ano & Por que ler \\
\midrule \endhead
Linguagem & \citetitle{livro:deitel2017}, 11ª ed. \pt{} \emph{Java: Como
  Programar}, 10ª ed. & \citeauthor{livro:deitel2017},
  \citeyear{livro:deitel2017} & Livro-texto para revisar a linguagem de
  ponta a ponta \\
Boas práticas & \citetitle{livro:bloch2017}, 3ª ed. \pt{} \emph{Java
  Efetivo} & \citeauthor{livro:bloch2017}, \citeyear{livro:bloch2017} &
  Itens 17 (imutabilidade) e 42 (lambdas); referência para decidir como
  escrever Java \\
Concorrência & \citetitle{livro:goetz2006} & \citeauthor{livro:goetz2006},
  \citeyear{livro:goetz2006} & Base de concorrência que continua valendo
  com virtual threads \\
\bottomrule
\end{longtable}
\normalsize

As edições citadas do Deitel e do \emph{Effective Java} são anteriores a
records, sealed types e virtual threads; para esses recursos, a fonte
primária são os JEPs da tabela do início do capítulo.

\subsection*{Ferramentas e documentação}
\citetitle{doc:oracle-virtual-threads}, \citetitle{doc:jls25} e os JEPs
citados. Para conferir o bytecode: \citetitle{doc:javap}.
